\scanpage{30}
\[
H^r(A,G)=(0)\quad\uline{\text{for}}\quad r>2.
\]
Proof. The duality theorem of [26] shows that $H^r(k,G')=(0)$ for $r>2$ if $G'$ is any AGS over $k$. Prop. 14 yields
\[
H^r(A,G)=(0)\quad\text{for}\quad r>3.
\]
To treat the case $r=3$, Th. 1 shows that we may assume $A$ is the field $k$. By Th. 4 and the duality theorem used above,
\[
H^3(k,G)\cong H^3(k,G_s).
\]
However, the right hand group vanishes by results of Tate [27], and the proof is complete.

One may obtain additional information concerning the cohomology groups with coefficients in $\underset{\sim}{G}_m$ by combining Th. 6 with results of Berkson, and Rosenberg-Zelinsky [5, 21]. Let $k$ be a field, and let $K$ be a purely inseparable extension of $k$. Let $\bigotimes^rK$ denote the $k$-algebra obtained by forming the tensor product of $K$ with itself (over $k$) $r$-times. The face operators
\[
\bigotimes^rK\to\bigotimes^{r+1}K
\]
induce group homomorphisms
\[
H^s(\bigotimes^rK,\underset{\sim}{G}_m)\to H^s(\bigotimes^{r+1}K,\underset{\sim}{G}_m).
\]
By combining the face operators in the obvious fashion, we obtain a complex, $C(s)$ for each $s>0$.
\[
C(s):H^s(K,\underset{\sim}{G}_m)\to H^s(K\otimes K,\underset{\sim}{G}_m)\to H^s(K\otimes K\otimes K,\underset{\sim}{G}_m)\to\cdots\ .
\]
Proposition 15. \uline{For each} $s>0$, \uline{the complex} $C(s)$ \uline{is acyclic}.

\begin{center}42\end{center}
\mspage{29}
